\section*{Hoofdstuk \ref{ch:b2:speciation} --- Soortvorming en macro-evolutie}

\begin{solution}{exo:b2:speciation:1}
Een soort is een groep natuurlijke populaties die onderling voortplanten en
reproductief geïsoleerd zijn van andere zulke groepen. Het begrip faalt voor
ongeslachtelijke lijnen (bacteriën, bdelloïde raderdiertjes, paardenbloemen:
morfologisch of fylogenetisch begrip, clusters van gelijke vormen of
herkenbare lijnen); voor fossielen (morfologisch begrip, omdat er geen
kruising kan worden gemaakt); en voor allopatrische populaties die elkaar nooit
ontmoeten (ecologisch of fylogenetisch begrip, omdat kruisbaarheid niet
kan worden getest) --- en het vervaagt door hybridiserende eiken en door
ringsoorten.
\end{solution}

\begin{solution}{exo:b2:speciation:2}
Kikkers in verschillende maanden: prezygotisch (temporeel). Muilezel:
postzygotisch (hybride steriliteit). Zaadcel die zich niet kan binden:
prezygotisch (gametisch). Krachtige hybriden met steriel \omterm{def:b2:angiosperm-reproduction:seed}{zaad}: postzygotisch
(hybride steriliteit, of hybride afbraak als het falen in de volgende
generatie ligt). Krekels met een verschillende zang: prezygotisch (gedrag).
\end{solution}

\begin{solution}{exo:b2:speciation:3}
Allopatrisch: een barrière verdeelt een verspreidingsgebied (pistoolgarnalen
aan de twee kanten van de landengte van Panama). Peripatrisch: enkele stichters
buiten het verspreidingsgebied (de vinken van Darwin uit een voorouder van het
vasteland). Sympatrisch: geen barrière (de appelboorvlieg op meidoorn en appel;
een allopolyploïde tarwe). Parapatrisch: langs een gradiënt met een
\omterm{prop:b2:speciation:reinforcement}{hybridisatiezone} in het midden (grassen op en naast mijnafval, die op
verschillende tijden bloeien).
\end{solution}

\begin{solution}{exo:b2:speciation:4}
Soort 1 wint altijd; soort 2 wint altijd; stabiele coëxistentie;
elk van beide wint afhankelijk van de beginaantallen. Coëxistentie is stabiel
wanneer elke soort de andere op haar draagkracht kan binnendringen,
dat wil zeggen wanneer elke soort zichzelf meer beperkt dan ze de
andere beperkt ($K_1 > \alpha K_2$ en $K_2 > \beta K_1$).
\end{solution}

\begin{solution}{exo:b2:speciation:5}
$\alpha K_2 = 400 < K_1 = 1000$: soort 1 dringt binnen; $\beta K_1 = 600
< K_2 = 800$: soort 2 dringt binnen. $N_1^{*} = (1000 - 400)/(1 - 0.3) =
857$, $N_2^{*} = (800 - 600)/0.7 = 286$; in het evenwicht ligt soort
1 onder haar $K$, met $\alpha N_2^{*} = 143$, en soort 2 met $\beta
N_1^{*} = 514$.
\end{solution}

\begin{solution}{exo:b2:speciation:6}
$\alpha K_2 = 1200 > K_1$: soort 1 kan soort 2 op haar draagkracht
niet binnendringen; $\beta K_1 = 300 < K_2$: soort 2 kan soort 1
binnendringen. Soort 2 wint altijd, omdat haar individuen soort 1 meer
schaden dan de eigen individuen van soort 1 dat doen, terwijl ze zelf
nauwelijks wordt getroffen. Met $\alpha = \beta = 1.5$: $\alpha K_2 = 1200 > 1000$
en $\beta K_1 = 1500 > 800$: geen van beide kan de andere binnendringen, elk
sluit de andere uit wanneer ze gevestigd is, en wie er eerst was, wint.
\end{solution}

\begin{solution}{exo:b2:speciation:7}
De tetraploïde $4n$ maakt gameten $2n$; gekruist met de gameten $n$ van de
diploïde geeft hij triploïden $3n$, waarvan de \omterm{def:b2:meiosis-heredity:meiosis}{meiose} geen drie
sets kan paren en aneuploïde, meestal niet levensvatbare gameten oplevert. De
tetraploïde kan alleen met zichzelf voortplanten (of door zelfbestuiving):
in één stap geïsoleerd. Hij faalt meestal omdat hij alleen is --- zijn enige
mogelijke partners zijn \omterm{def:b2:meiosis-heredity:meiosis}{diploïden} en elke zulke kruising is verspild ---
en wordt overspoeld; hij blijft bestaan wanneer hij zichzelf bestuift, zich
vegetatief vermeerdert, of herhaaldelijk ontstaat.
\end{solution}

\begin{solution}{exo:b2:speciation:8}
$0.01\times 10\times 10 = 1$ onverenigbaarheid; met elk 30, $0.01\times
900 = 9$. Isolatie groeit als het kwadraat van de divergentietijd: twee
populaties die drie keer zo lang uiteengelopen zijn, zijn negen keer zo
geïsoleerd --- de sneeuwbal.
\end{solution}

\begin{solution}{exo:b2:speciation:9}
$w = 2\sqrt{8/0.05} = \qty{25}{km}$; $w = 0.3\sqrt{8/0.4} =
\qty{1.3}{km}$. Versterking is waarschijnlijker in de tweede: hybriden
zijn sterk ongeschikt, zodat een vrouwtje dat de andere vorm weigert
veel wint, en de zone is smal genoeg om de twee vormen vaak genoeg te laten
ontmoeten om die weigering te laten selecteren.
\end{solution}

\begin{solution}{exo:b2:speciation:10}
$\alpha K_2 = 450 < 500$, dus elk dringt de andere binnen: ze leven samen,
bij $N^{*} = (500 - 450)/(1 - 0.81) = 263$ elk --- nipt, dicht bij de
grens van uitsluiting, en een kleine verandering van $K$ zou het doen omslaan.
Met $\alpha = \beta = 0.4$: $N^{*} = (500 - 200)/(1 - 0.16) = 357$ elk.
De verschuiving draaide elke isocline weg van die van de andere ($K/\alpha$
schoof naar buiten), zodat het snijpunt weg van de assen schoof en elke
soort dichter bij haar eigen $K$ zit: een zekerdere coëxistentie met
meer van elk.
\end{solution}

\begin{solution}{exo:b2:speciation:11}
$R = h^{2}S = 0.7\times(-1) = \qty{-0.7}{mm}$: zoals waargenomen. Bij de
droogte van 1977, zonder grote grondvink, werden dezelfde middelste
grondvinken geselecteerd op \emph{grotere} snavels, omdat grote zaden het
voedsel waren dat overbleef; in 2004 werden de grote zaden door de
indringer ingenomen, en verloren de middelste grondvinken die erom
concurreerden: de richting van de selectie werd door de concurrent bepaald.
\end{solution}

\begin{solution}{exo:b2:speciation:12}
Bij toeval ontstaan: het model met twee loci laat zien dat isolatie ontstaat
als bijproduct van substituties die om andere redenen of door drift zijn
gefixeerd, en de geografie --- een barrière, een stichting --- levert de
scheiding die het toeval zich laat opstapelen. Door selectie in stand
gehouden: waar de vormen elkaar ontmoeten, bouwt selectie tegen ongeschikte
hybriden prezygotische barrières op (versterking), en verschuift concurrentie
kenmerken zodat de twee kunnen samenleven. De slogan laat \omterm{prop:b2:speciation:modes}{sympatrische
soortvorming} door disruptieve selectie weg, waar selectie de soort zowel maakt
als in stand houdt, en \omterm{def:b2:genome-diversification:polyploidy}{polyploïdie}, waar het toeval het hele verhaal is.
\end{solution}

\begin{solution}{pb:b2:speciation:1}
\textbf{1.} Soort 1: $N_1 + 0.9N_2 = 1200$, snijpunten $N_1 = 1200$
en $N_2 = 1333$. Soort 2: $N_2 + 0.3N_1 = 400$, snijpunten $N_2 =
400$ en $N_1 = 1333$.
\textbf{2.} $\beta K_1 = 360 < K_2 = 400$: soort 2 dringt binnen;
$\alpha K_2 = 360 < K_1 = 1200$: soort 1 dringt binnen. Stabiele
coëxistentie.
\textbf{3.} $N_1^{*} = (1200 - 360)/(1 - 0.27) = 1151$; $N_2^{*} =
(400 - 360)/0.73 = 55$.
\textbf{4.} $\mathrm{d}N_1/\mathrm{d}t = 0.5\times 1200\,(1 - 1218/1200)
= -9$ per jaar: de inheemse soort, op haar draagkracht, wordt door de
nieuwkomers in een daling geduwd.
\textbf{5.} $K_1 = 600$, $K_2 = 200$: $N_1^{*} = (600 - 180)/0.73 =
575$, $N_2^{*} = (200 - 180)/0.73 = 27$: beide halveren, en de indringer
staat met 27 vogels dicht bij uitsterven --- in de droogte bijt de
concurrentie.
\textbf{6.} De indringer is de grotere vogel met de grotere snavel: hij
neemt de grote zaden die de inheemse soort ook gebruikt, en neemt ze
sneller, terwijl de kleine zaden van de inheemse soort voor hem weinig nut
hebben.
\textbf{7.} $S = 9.0 - 9.5 = \qty{-0.5}{mm}$; $R = 0.7\times(-0.5) =
\qty{-0.35}{mm}$.
\textbf{8.} $9.5 - 0.35 = \qty{9.15}{mm}$.
\textbf{9.} $N_1^{*} = (1200 - 200)/(1 - 0.15) = 1176$; $N_2^{*} =
(400 - 360)/0.85 = 47$.
\textbf{10.} Het snijpunt van de isocline van soort 1 met de $N_2$-as steeg van
$K_1/0.9 = 1333$ tot $K_1/0.5 = 2400$: de lijn draaide weg van die
van de indringer, de inheemse soort wordt nu minder getroffen door elke
indringer, en het evenwicht schoof naar de $K$ van de inheemse soort.
\textbf{11.} $2/0.35 \approx 6$ jaar. Dat gebeurt niet, omdat selectie
van deze soort alleen in droge jaren optreedt, in natte jaren omkeert
wanneer kleine zaden overvloedig zijn (of wanneer de indringer schaars is), en
omdat de additieve variantie uitgeput zou raken.
\textbf{12.} Twee miljoen jaar is ongeveer $10^{6}$ generaties; een
verschuiving van \qty{0.35}{mm} per generatie die zelfs een duizendste
daarvan aanhield, zou de snavel met honderden millimeters veranderen. Selectie
is sterk maar wisselvallig; aanhoudende gerichte verandering is de zeldzame
uitzondering, en de netto verandering over geologische tijd is een klein
restant van grote tegengestelde episodes.
\textbf{13.} De zang, die van de vader wordt geleerd en door vrouwtjes bij de
partnerkeuze wordt gebruikt, met de snavelvorm zelf als signaal. Hybriden met
tussenliggende snavels verwerken noch de grote noch de kleine zaden het best,
en verhongeren als eerste bij droogte.
\textbf{14.} $w = 0.5\sqrt{8/0.2} = \qty{3.2}{km}$.
\textbf{15.} Het eiland is smaller dan de zone zou zijn: er kan zich geen
ruimtelijke cline vormen; hybridisatie treft, waar ze optreedt, de
hele populatie, en alleen assortatieve paring houdt de twee
gescheiden.
\textbf{16.} Vrouwtjes die alleen paren met mannetjes met hun eigen zang en
snavel, laten meer kleinkinderen na dan vrouwtjes die hybridiseren; de
voorkeur en het signaal (zang, snavel) worden duidelijker verschillend waar de
twee soorten samenleven dan waar ze gescheiden leven.
\textbf{17.} $0.005\times 15\times 15 = 1.1$ onverenigbaarheden: gemiddeld
één, met hybriden die waarschijnlijk minder fit maar niet steriel zijn. Volgens
het biologische criterium nog geen soorten; wel op weg.
\textbf{18.} Biologisch: nog niet (hybriden bijna fit).
Morfologisch: misschien, als de snavels of het verenkleed zichtbaar
uiteengelopen zijn. Fylogenetisch: ja, als elke populatie een gefixeerd
diagnostisch verschil heeft. De begrippen zijn het oneens precies tijdens het
proces van soortvorming, en dat maakt ze nuttig.
\textbf{19.} $1/(5\times 10^{6}) = 2\times 10^{-7}$ per soort per
jaar; voor $10^{7}$ soorten twee uitstervingen per jaar.
\textbf{20.} 200 tot 2000 soorten per jaar.
\textbf{21.} $0.75\times 10^{7}/10^{5} = 75$ soorten per jaar --- lager
dan de huidige snelheid aan de onderkant van de schatting. De huidige episode
is, als ze aanhoudt, een massa-extinctie die sneller verloopt dan de grote
vijf.
\textbf{22.} De niches die de verliezers achterlaten, zijn vrij van
concurrenten, zodat elke overlevende variant die er een kan benutten wordt
begunstigd, en de overlevenden zich erin diversifiëren.
\textbf{23.} \omterm{prop:b2:speciation:tempo}{Onderbroken evenwicht}; allopatrische (peripatrische)
soortvorming in een kleine geïsoleerde populatie waarvan de afstammeling zich
daarna over het verspreidingsgebied van de hoofdpopulatie uitbreidt.
\textbf{24.} $1.04^{n} = 2$: $n = \ln 2/\ln 1.04 = 18$ generaties.
Het fossielenbestand toont miljoenen jaren omdat selectie omkeert,
\omterm{prop:b2:population-genetics:kinds}{stabiliserende selectie} meestal overheerst, en de netto verandering het kleine
verschil is tussen grote tegengestelde episodes.
\textbf{25.} Vóór de verschuiving: stabiele coëxistentie bij $N_1^{*} =
1151$, $N_2^{*} = 55$; erna: 1176 en 47. Verschuiving van de snavel
\qty{-0.35}{mm} in één generatie. Breedte van de \omterm{prop:b2:speciation:reinforcement}{hybridisatiezone} \qty{3.2}{km}.
\end{solution}
